\section{Introduction}
\label{sec:intro}

A prediction market turns a question about the future into a contract that pays
one dollar if the answer is Yes and nothing if it is No. Its trading price is a
probability with money behind it, and decades of work have found those prices
hard to beat~\citep{wolfers2004}. Recent studies report that large language
models now forecast about as well as human crowds~\citep{schoenegger2024}, beat
the median forecaster on public tournaments~\citep{lu2025}, and, in one
recent workshop paper on prediction markets, match the price on accuracy while
earning higher returns~\citep{henry2026beyond}. If that is right, it is a notable claim about
machine judgment. It is also hard to test, because three things have to be true
of the comparison.

\textbf{The market has to be trading.} On both venues a fresh listing carries a
one-cent bid against a ninety-nine-cent ask, and the midpoint is fifty cents whatever the event's real chance. Every row of the two live-book windows has a live book at the
forecast moment: quotes on both sides, with a spread of ten cents or less. The third, \polybench's published week, was drawn without that rule, and we rest
no claim on it alone.

\textbf{The model has to form its own number.} If a model is shown the market
price, whatever it returns is contaminated by the number it was handed. Language
models move markedly toward a crowd estimate when shown
one~\citep{schoenegger2024}. No model in this study sees the price, the best bid,
the best ask, the depth or the book; it sees the question, the resolution rules,
the scheduled close and a frozen news snapshot.

\textbf{The model must not already know the answer.} Restricting retrieval by date
does not prevent leakage: one audit found major post-cutoff information in 71
and 81 per cent of date-filtered pages, with the answer revealed outright for 41
and 55 per cent~\citep{ellahib2026}. A model forecasts a row only if its exact published
version was released before that row's outcome could be known, the news was
captured before the forecast moment, and no model performs a live web search.

With those three conditions in place, the market is more accurate than every
model in every depth bucket on the two live-book windows. The gap is resolution
rather than calibration: no monotone rescaling recovers more than about a fifth
overall. The losses concentrate where a model contradicts the crowd
(Section~\ref{sec:results-crowd}). The models also err together: on rows the
market got right, two models land on the same wrong side 2.7 to 4.3 times as
often as independent forecasters would. Controlling for what they were jointly
given leaves that rate well above what those inputs alone produce, and adding
news lowers their agreement only slightly; information none of them had, shared training and a shared prompt
remain. No combination of the models reaches the market
(Section~\ref{sec:results-correlated}).

We contribute a replay benchmark of 10{,}536 market moments on two venues, a
price-blind evaluation harness (the code that prompts each model and records its
answer) run over eight release-date-eligible models, the finding
that the accuracy gap lies on the resolution axis, checked against recalibration and
learned combination, and a correlated-error measure tested against shared
shrinkage, the observable features of a row and the event itself.

\begin{figure*}[!tbp]
\centering
\begin{tikzpicture}
\begin{groupplot}[
  paperaxis,
  group style={group size=2 by 1, horizontal sep=6pt},
  width=0.5\linewidth, height=0.35\linewidth,
  xmin=0, xmax=1, ymin=0, ymax=1,
  xtick={0,0.25,0.5,0.75,1}, ytick={0,0.25,0.5,0.75,1},
  xlabel={Market price (mean in band)},
  grid=major,
  legend columns=3,
]
\nextgroupplot[title={\winB (\polymarket)}, ylabel={Mean model forecast}, legend to name=legFc,]
\addplot[mdl/market, forget plot, no marks, line width=0.6pt] coordinates {(0,0) (1,1)};
\addplot[black, forget plot, solid, line width=1.0pt, mark=*, mark size=1.6pt] coordinates {(0.134,0.218) (0.391,0.397) (0.605,0.514) (0.852,0.604)};
\addplot[black, forget plot, densely dotted, line width=1.0pt, mark=o, mark size=1.6pt, mark options={solid}] coordinates {(0.134,0.201) (0.391,0.380) (0.605,0.470) (0.852,0.507)};
\addplot[gray, forget plot, only marks, mark=x, mark size=2.6pt, line width=0.8pt] coordinates {(0.134,0.128) (0.391,0.346) (0.605,0.585) (0.852,0.907)};
\addplot[mdl/gpt-5.2, forget plot, only marks, papermark] coordinates {(0.148,0.207) (0.846,0.579)};
\addplot[mdl/gpt-oss-120b, forget plot, only marks, papermark] coordinates {(0.148,0.228) (0.846,0.490)};
\addplot[mdl/kimi-k2.5, forget plot, only marks, papermark] coordinates {(0.148,0.220) (0.846,0.563)};
\addplot[mdl/deepseek-v3.2, forget plot, only marks, papermark] coordinates {(0.148,0.245) (0.846,0.575)};
\addplot[mdl/gpt-5.5, forget plot, only marks, papermark] coordinates {(0.148,0.209) (0.846,0.607)};
\addplot[mdl/glm-5.1, forget plot, only marks, papermark] coordinates {(0.148,0.225) (0.846,0.608)};
\addplot[mdl/kimi-k2.6, forget plot, only marks, papermark] coordinates {(0.147,0.216) (0.846,0.582)};
\addplot[mdl/deepseek-v4-pro, forget plot, only marks, papermark] coordinates {(0.148,0.215) (0.846,0.637)};
\addlegendimage{mdl/gpt-5.2, only marks, papermark}\addlegendentry{\mGptFiveTwo}
\addlegendimage{mdl/gpt-oss-120b, only marks, papermark}\addlegendentry{\mGptOss}
\addlegendimage{mdl/kimi-k2.5, only marks, papermark}\addlegendentry{\mKimiKTwoFive}
\addlegendimage{mdl/deepseek-v3.2, only marks, papermark}\addlegendentry{\mDeepSeekVThreeTwo}
\addlegendimage{mdl/gpt-5.5, only marks, papermark}\addlegendentry{\mGptFiveFive}
\addlegendimage{mdl/glm-5.1, only marks, papermark}\addlegendentry{\mGlmFiveOne}
\addlegendimage{mdl/kimi-k2.6, only marks, papermark}\addlegendentry{\mKimiKTwoSix}
\addlegendimage{mdl/deepseek-v4-pro, only marks, papermark}\addlegendentry{\mDeepSeekVFourPro}
\addlegendimage{mdl/market, no marks, line width=0.6pt}\addlegendentry{Market price (diagonal)}
\addlegendimage{black, solid, line width=1.0pt, mark=*, mark size=1.6pt}\addlegendentry{All models, with news}
\addlegendimage{black, densely dotted, line width=1.0pt, mark=o, mark size=1.6pt, mark options={solid}}\addlegendentry{All models, news removed}
\addlegendimage{gray, only marks, mark=x, mark size=2.6pt, line width=0.8pt}\addlegendentry{Share resolved Yes}
\nextgroupplot[title={\winC (\kalshi)}, yticklabels={},]
\addplot[mdl/market, forget plot, no marks, line width=0.6pt] coordinates {(0,0) (1,1)};
\addplot[black, forget plot, solid, line width=1.0pt, mark=*, mark size=1.6pt] coordinates {(0.129,0.188) (0.366,0.400) (0.620,0.545) (0.846,0.629)};
\addplot[black, forget plot, densely dotted, line width=1.0pt, mark=o, mark size=1.6pt, mark options={solid}] coordinates {(0.129,0.166) (0.366,0.370) (0.620,0.519) (0.846,0.574)};
\addplot[gray, forget plot, only marks, mark=x, mark size=2.6pt, line width=0.8pt] coordinates {(0.129,0.142) (0.366,0.370) (0.620,0.635) (0.846,0.867)};
\addplot[mdl/gpt-5.2, forget plot, only marks, papermark] coordinates {(0.132,0.203) (0.843,0.640)};
\addplot[mdl/gpt-oss-120b, forget plot, only marks, papermark] coordinates {(0.132,0.222) (0.843,0.546)};
\addplot[mdl/kimi-k2.5, forget plot, only marks, papermark] coordinates {(0.132,0.211) (0.843,0.601)};
\addplot[mdl/deepseek-v3.2, forget plot, only marks, papermark] coordinates {(0.132,0.235) (0.843,0.587)};
\addplot[mdl/gpt-5.5, forget plot, only marks, papermark] coordinates {(0.132,0.196) (0.843,0.642)};
\addplot[mdl/glm-5.1, forget plot, only marks, papermark] coordinates {(0.132,0.221) (0.843,0.624)};
\addplot[mdl/kimi-k2.6, forget plot, only marks, papermark] coordinates {(0.132,0.210) (0.843,0.598)};
\addplot[mdl/deepseek-v4-pro, forget plot, only marks, papermark] coordinates {(0.132,0.205) (0.843,0.670)};

\end{groupplot}
\end{tikzpicture}

\pgfplotslegendfromname{legFc}

\caption{No model sees the market price, yet the forecasts flatten against it. Black lines
are all models pooled on news-bearing rows, with (solid) and without (dotted) the news, in
four price bands. Grey crosses are the actual Yes share on those rows. Coloured markers are each model's
mean forecast in the lowest and highest band, over all primary rows. Test split; pooled lines
use news-bearing rows, per-model markers use every row. Top-band Yes share: 91 and 87 per cent
on news-bearing rows, 86 and 84 on all.}
\label{fig:forecast-vs-price}
\end{figure*}
